\section{Interpretation, limitations and next evidence}
The strongest supported result is a separation of alarm entry quality from evidence for ending an alarm. At fixed entry threshold, a dwell and reference-distance condition prevented observed premature clears, while increasing benign alarm time. The command-interval ranking favored ordinary following hysteresis. A claim of overall improvement would hide one of these outcomes. The tested mechanism is useful as a diagnostic design and a controlled counterexample to treating F1 or alarm persistence alone as sufficient evidence of good recovery behavior.

This result is not a demonstrated flaw in CASPER or its successor. Those papers use different sensing modalities and their published experiments have not been faithfully reproduced here. The present classical controls were implemented explicitly to test a mechanism under common information. The earlier IMU experiments remain supporting development records, not independent confirmation of these simulated cyber conclusions. A recent-paper shortlist cannot substitute for compatible input access and a traceable experiment.

The return diagnostic is itself a substantive assumption. True state is compared with a nominal reference, not a physically measured safe set or a production-quality criterion. Recovery could require a different tolerance or task-specific condition; the present threshold was frozen rather than retrospectively justified. The fixed reference also continues to suffer model mismatch. Gaussian noise is unbounded, the sampled envelope has no certified containment property, and higher fresh noise violates the nominal margin assumption. A safety guarantee cannot be inferred from the absence of premature clears in 90 events.

The evidence spans two new analytic tasks and three plant settings but one modeling family. Its channels, controller and disturbance forms are small compared with an industrial robot. Paired seeds change sensing noise on the same deterministic physical trajectories. They do not represent independent tasks, installations or attack campaigns. Parameters were not estimated from a newly admitted real plant, and no authenticated command channel was implemented. The new task/lag challenge is meaningful within this declared model, while cross-simulator or real-recording validity remains open.

Future work can remain fully virtual. A defensible next mechanism would need a separate, causally available indication that the model is adequate during legitimate changes, without using true test parameters or attack labels. Before implementation, it must be distinguished from interval-observer, residual-consistency and recovery methods in the closest literature. Its proof obligations should specify what prevents unsafe evidence removal and what assumptions bound the added benign burden. An independent implementation of different dynamics, with newly frozen task/plant conditions, would strengthen that test. Reusing this now-exposed condition matrix after redesign would be exploratory.

Public recordings could add a separate validity track. The voraus-AD author resources provide accessible robotic fault data, including measured and target signals. A target signal is not automatically an authenticated approved command, and physical faults are not automatically cyberattacks. Admission should first establish acquisition units, release hashes, feature trust and a question answerable by their labels. A fault-generalization result may be valuable in its own right, but it must not be merged with the present cyber-recovery denominator. No physical hardware is required to pursue either the independent simulation or public-data route.

\section{Conclusion}
A separately defined fixed-return rule changed recovery behavior in a frozen virtual motion-monitoring study. Its dwell prevented premature clearing observed in 44 of 90 matched-threshold no-dwell events, while reference-based entry added no alarm-decision benefit in the tested matrix. The candidate nevertheless increased benign alarm points by \RelativeIncrease\% relative to fused hysteresis and failed the predeclared overall improvement target. The current evidence supports reporting recovery, fresh detection and benign burden together. It does not establish a new algorithm, superior baseline performance or a physical safety guarantee. An original contribution now requires a specific mechanism that resolves this measured trade-off under fresh evidence, rather than a new name for the existing latch.
